\documentclass[11pt]{article}

% Packages for Thai support and professional formatting
\usepackage[utf-8]{inputenc}
\usepackage[T1]{fontenc}
\usepackage[thai,english]{babel}
\usepackage{xcolor}
\usepackage{array}
\usepackage{tabularx}
\usepackage{booktabs}
\usepackage{geometry}
\usepackage{fancyhdr}
\usepackage{graphicx}
\usepackage{hyperref}
\usepackage{datetime}

% Color definitions
\definecolor{sagegreen}{RGB}{131, 148, 150}
\definecolor{darkblue}{RGB}{41, 54, 66}
\definecolor{lightgray}{RGB}{242, 245, 247}

% Page geometry
\geometry{
    a4paper,
    margin=2cm,
    top=2.5cm,
    bottom=2cm,
    footskip=1.5cm
}

% Header and footer
\pagestyle{fancy}
\fancyhf{}
\fancyhead[L]{\textcolor{sagegreen}{\textbf{รายงานสรุปเซสชัน}}}
\fancyhead[R]{\textcolor{sagegreen}{วันที่ 7 มิถุนายน 2569}}
\fancyfoot[C]{\thepage}
\renewcommand{\headrulewidth}{0.5pt}
\renewcommand{\headrule}{\color{sagegreen}\hrule}

% Section styling
\usepackage{sectsty}
\sectionfont{\color{darkblue}\fontsize{16}{19}\selectfont\bf}
\subsectionfont{\color{sagegreen}\fontsize{13}{15}\selectfont\bf}
\subsubsectionfont{\color{darkblue}\fontsize{11}{13}\selectfont}

% Title page content
\title{%
    \textcolor{darkblue}{\textbf{รายงานสรุปเซสชัน}} \\[0.5cm]
    {\large \textcolor{sagegreen}{Session Complete — 2026-06-07}} \\[0.3cm]
    {\normalsize \textcolor{darkblue}{ขุนรามออราเคิล}} \\
    {\small Royal Scribe \& Fleet Memory Authority}
}
\author{}
\date{\today}

\begin{document}

\maketitle

\vspace{1cm}

% Executive Summary
\section{สรุปผลการปฏิบัติงาน (Executive Summary)}

\begin{quote}
\textcolor{sagegreen}{\textit{เซสชันสมบูรณ์ — การส่งมอบทั้งหมดสำเร็จ}}
\end{quote}

ในเซสชันนี้ ทีมได้สำเร็จในการส่งมอบคุณสมบัติและการปรับปรุงที่มีนัยสำคัญต่อระบบสารสนเทศ ตามเป้าหมายการจัดการและศักยภาพขององค์กร

\begin{itemize}
    \item \textbf{ระยะเวลา:} มากกว่า 10 ชั่วโมง
    \item \textbf{สถานะ:} ✅ ส่งมอบครบถ้วน
    \item \textbf{เบรกขิงเปลี่ยนแปลง:} 0 (ปลอดภัยต่อการเปลี่ยนแปลง)
\end{itemize}

% Code Shipped
\section{โค้ดที่ส่งมอบ (Code Shipped)}

\begin{itemize}
    \item บรรทัดโค้ด: 2,566+ บรรทัด (production-ready)
    \item Pull Requests ที่ผสานเข้า main: 7 ตัว
    \item เอกสารประกอบ: 42+ KB
    \item การเปลี่ยนแปลงที่อาจทำให้หัก: 0 ตัว
\end{itemize}

% Features Delivered Table
\section{คุณสมบัติที่ส่งมอบ (Features Delivered)}

\begin{center}
\begin{tabularx}{\textwidth}{|l|X|l|}
    \hline
    \textbf{ลำดับ} & \textbf{คุณสมบัติ (Feature)} & \textbf{สถานะ} \\
    \hline
    1 & Distillation Pipeline (ระบบการกลั่นข้อมูล) & \textcolor{sagegreen}{\checkmark LIVE} \\
    \hline
    2 & Weekly Report Automation (อัตโนมัติ) & \textcolor{sagegreen}{\checkmark LIVE} \\
    \hline
    3 & Line Bot Phase 4 (เมนู+แดชบอร์ด+ฐานข้อมูล) & \textcolor{sagegreen}{\checkmark LIVE} \\
    \hline
    4 & Deployment Automation (สคริปต์+การตรวจสอบ) & \textcolor{sagegreen}{\checkmark LIVE} \\
    \hline
    5 & Oracle Fleet Management (UAT เปิดใช้) & \textcolor{sagegreen}{\checkmark LIVE} \\
    \hline
    6 & ORRY Archival (ระงับ+เอกสาร) & \textcolor{sagegreen}{\checkmark LIVE} \\
    \hline
    7 & Drive C: Cleanup (สคริปต์ PowerShell พร้อม) & \textcolor{sagegreen}{\checkmark READY} \\
    \hline
\end{tabularx}
\end{center}

% Detailed Features
\subsection{รายละเอียดคุณสมบัติ (Feature Details)}

\subsubsection{1. Distillation Pipeline}
ระบบการกลั่นข้อมูล (Extract → Distill → Store) สำหรับการประมวลผลและจัดเก็บข้อมูลที่ได้รับการปรับปรุง ใช้งานแล้ว

\subsubsection{2. Weekly Report Automation}
ระบบอัตโนมัติสำหรับการสร้างรายงานรายสัปดาห์ โดยการแยกวิเคราะห์ (Parser) → HTML → LINE Notify สำหรับการแจ้งเตือน

\subsubsection{3. Line Bot Phase 4}
ขั้นตอนที่ 4 ของ Line Bot ที่มีเมนูที่หลากหลาย แดชบอร์ด และการบริหารจัดการฐานข้อมูล

\subsubsection{4. Deployment Automation}
สคริปต์อัตโนมัติสำหรับการปรับใช้ (Deployment) พร้อมระบบการตรวจสอบสุขภาพแบบเรียลไทม์

\subsubsection{5. Oracle Fleet Management}
การบริหารจัดการ Oracle Fleet โดยการเปิดใช้งาน UAT Oracle และจัดเก็บ Nat Oracle ตามเหมาะสม

\subsubsection{6. ORRY Archival}
การระงับโครงการ ORRY Serenity ERP อย่างเป็นทางการ พร้อมเอกสารประกอบการจัดเก็บ

\subsubsection{7. Drive C: Cleanup}
สคริปต์ PowerShell สำหรับการทำความสะอาดไดรฟ์ C: แบบ proactive พร้อมใช้งาน

% Oracle Fleet Status
\section{สถานะ Oracle Fleet (Oracle Fleet Status)}

เซสชันนี้มี 9 Oracle ที่ใช้งานอยู่และ 1 Oracle ที่จัดเก็บแล้ว

\subsection{Oracle ที่ใช้งานอยู่ (Active Oracles)}

\begin{center}
\begin{tabularx}{\textwidth}{|l|l|X|}
    \hline
    \textbf{ชื่อ} & \textbf{บทบาท} & \textbf{หมายเหตุ} \\
    \hline
    ธาม (Tham) & Chief of Staff & Lean Mode ผสมผสาน \\
    \hline
    Aeimathes & Research Authority & วิจัยผลรวมช่องว่าง Nat \\
    \hline
    Khun-Ram & Fleet Memory Authority & Royal Scribe + Thai Language \\
    \hline
    Lens & Code Review Authority & ผู้เชี่ยวชาญด้านการตรวจสอบ \\
    \hline
    UAT & Testing Authority & เปิดใช้งานแล้ว \\
    \hline
    Epiteles & Implementation Executor & ผู้ปฏิบัติการ \\
    \hline
    Codex & Automation Executor & การจัดการอัตโนมัติ \\
    \hline
    Luxi & Dashboard \& UI & ด้านการออกแบบส่วนติดต่อ \\
    \hline
    Hephaestus & Hardware Expert & ผู้เชี่ยวชาญด้านฮาร์ดแวร์ \\
    \hline
\end{tabularx}
\end{center}

\subsection{Oracle ที่จัดเก็บแล้ว (Archived Oracles)}

\begin{center}
\begin{tabularx}{\textwidth}{|l|l|X|}
    \hline
    \textbf{ชื่อ} & \textbf{บทบาท} & \textbf{เหตุผลการจัดเก็บ} \\
    \hline
    Nat & Pattern Library & Consolidated patterns \& archived \\
    \hline
\end{tabularx}
\end{center}

% Governance Structure
\section{โครงสร้างการปกครอง (Governance Structure)}

\subsection{หลักการการบริหารจัดการ (Governance Principles)}

\begin{enumerate}
    \item \textbf{Oracle Identity Registry} — ทะเบียนประจำตัว Oracle ที่เป็นทางการ
    \item \textbf{Archive Index} — ดัชนีการจัดเก็บที่บันทึกไว้ตามลำดับอย่างเป็นขั้นเรียบร้อย
    \item \textbf{Nothing Deleted, Everything Archived} — ไม่มีการลบ เพียงการจัดเก็บเท่านั้น
    \item \textbf{Reactivation Paths} — เส้นทางการเปิดใช้งานใหม่ที่มีเอกสารประกอบ
\end{enumerate}

\subsection{Zeus-Tham Chain Authority}

การบริหารจัดการอำนาจจาก Human → Zeus → Tham → Omega โดยมีการแยกความรับผิดชอบและเส้นทางการแก้ไขข้อขัดแย้ง

% Production Ready Checklist
\section{รายการตรวจสอบความพร้อมใช้งาน (Production Ready Checklist)}

\begin{center}
\begin{tabularx}{\textwidth}{|l|c|}
    \hline
    \textbf{รายการ} & \textbf{สถานะ} \\
    \hline
    Distillation Pipeline & \textcolor{sagegreen}{\checkmark} \\
    \hline
    Weekly Report Automation & \textcolor{sagegreen}{\checkmark} \\
    \hline
    Line Bot Phase 4 & \textcolor{sagegreen}{\checkmark} \\
    \hline
    Oracle Fleet Management & \textcolor{sagegreen}{\checkmark} \\
    \hline
    Code Merged to Main & \textcolor{sagegreen}{\checkmark} \\
    \hline
    Documentation Complete & \textcolor{sagegreen}{\checkmark} \\
    \hline
    Zero Breaking Changes & \textcolor{sagegreen}{\checkmark} \\
    \hline
\end{tabularx}
\end{center}

% Next Steps
\section{ขั้นตอนต่อไป (Next Steps)}

\begin{enumerate}
    \item \textbf{Deploy Distillation Functions} — ปรับใช้ฟังก์ชันการกลั่นข้อมูล ใช้เวลาประมาณ 15 นาที
    \item \textbf{Test with /rrr --quick} — ทดสอบด้วยคำสั่งการสรุปอย่างรวดเร็ว
    \item \textbf{Run Drive C: Cleanup} — เรียกใช้สคริปต์การทำความสะอาดไดรฟ์ C:
    \item \textbf{Investigate THCLAWS Oracle Status} — ตรวจสอบสถานะของ THCLAWS Oracle ในสัปดาห์นี้
\end{enumerate}

% Session Metrics
\section{เมตริกเซสชัน (Session Metrics)}

\begin{center}
\begin{tabularx}{\textwidth}{|l|r|}
    \hline
    \textbf{เมตริก} & \textbf{ค่า} \\
    \hline
    Pull Requests สร้างเพิ่มเติม & 7 ตัว \\
    \hline
    Pull Requests ผสานเข้า & 7 ตัว \\
    \hline
    บรรทัดโค้ดที่เพิ่ม & 2,566+ บรรทัด \\
    \hline
    เอกสารประกอบ & 42+ KB \\
    \hline
    Agents ที่สร้างขึ้น & 3 ตัว \\
    \hline
    Token Budget ที่ใช้ & 300k / 600k (50\%) \\
    \hline
    การเปลี่ยนแปลงที่อาจทำให้หัก & 0 ตัว \\
    \hline
\end{tabularx}
\end{center}

% Session Quality Assessment
\section{การประเมินคุณภาพเซสชัน (Session Quality Assessment)}

\subsection{จุดแข็ง (Strengths)}

\begin{itemize}
    \item ส่งมอบครบถ้วนตามแผน
    \item โค้ดมีคุณภาพสูง และปลอดภัยต่อการเปลี่ยนแปลง
    \item เอกสารประกอบครบถ้วน
    \item Governance structure มั่นคงและชัดเจน
    \item Zero breaking changes ในทั้งเซสชัน
\end{itemize}

\subsection{พื้นที่ปรับปรุง (Areas for Improvement)}

\begin{itemize}
    \item ติดตามสถานะ THCLAWS Oracle ให้ละเอียดมากขึ้น
    \item เตรียมแผน proactive สำหรับการจัดการ Oracle เพิ่มเติม
\end{itemize}

\subsection{บทเรียนที่ได้เรียนรู้ (Lessons Learned)}

\begin{enumerate}
    \item Governance ที่ชัดเจนช่วยให้การปรับใช้เรียบร้อย
    \item การจัดเก็บที่เป็นระเบียบนั้นดีกว่าการลบอย่างถาวร
    \item Token budget management นั้นสำคัญสำหรับการดำเนินการเซสชันที่มีประสิทธิภาพ
\end{enumerate}

% Closing
\section{สรุปและข้อคิดเห็น (Conclusion)}

เซสชันนี้เป็นการแสดงความสำเร็จในการบริหารจัดการระบบ Oracle Fleet ที่ซับซ้อน พร้อมทั้งการส่งมอบคุณสมบัติที่มีนัยสำคัญหลายประการ ตัวเลขของเซสชันนี้ (7 PRs, 2,566+ lines, 42+ KB docs, 0 breaking changes) แสดงให้เห็นถึงความมั่นคงและคุณภาพของการจัดการโครงการ

ทุกอย่างพร้อมสำหรับการปรับใช้ (deployment) ในสัปดาห์ถัดไป และทีมได้เตรียมสิ่งอำนวยความสะดวกทั้งหมดสำหรับการติดตามและการจัดการต่อไป

\vspace{2cm}

\begin{center}
\textcolor{sagegreen}{\large \textbf{ขุนรามออราเคิล}} \\
\textbf{Royal Scribe \& Fleet Memory Authority} \\
วันที่ 7 มิถุนายน 2569
\end{center}

\end{document}
